% Standalone appendix document - compiles on its own:
%   latexmk -pdf -outdir=build 07-agenda-2026-11-03.tex
\documentclass[a4paper,11pt]{article}
\usepackage[danish]{babel}
\usepackage{../../appendix-style}

\title{12.7 Dagsorden -- 03-11}
\author{SW4PRJ4 -- Group 3}
\date{03-11-2026}

\begin{document}
\maketitle

\section*{Mødeinformation}

\textbf{Dato:} 03-11-2026 (tirsdag)

\textbf{Tid:} 13:00

\textbf{Sted:} 5125-110

\textbf{Deltagere:} Gruppe 3 og vejleder Michel

\section*{Dagsordenspunkter}

\begin{enumerate}
    \item \textbf{Status siden sidste vejledermøde (6/10)}
    \begin{itemize}
        \item I sprint 3 fik vi login med Keycloak i kunde-appen og admin-panelet. En ny kunde kan oprette sig med navn og adresse, og backend kender forskel på en kunde og en admin.
        \item Kunde-appen husker nu login sikkert på telefonen, og admin-panelet er deployet.
        \item Controllerne går nu gennem services og ikke direkte til databasen.
        \item Vi har tegnet diagrammer i draw.io over login med Keycloak, onboarding af kunde og admin, hentning af items, hvordan en kunde bliver gemt, pakkerne i koden, pipeline og rollback, sprint-forløbet og et rich picture. De fleste er sat ind i rapporten og bilagene. Til login, onboarding og hentning af items mangler vi stadig at skrive teksten.
        \item I rapporten har vi skrevet kapitel 1 og risikoanalysen. Vi har også gennemgået den tekniske analyse, kravspecifikationen og accepttesten.
        \item Vi lovede at have abonnementer og ordrer i databasen til i dag. Vi fortæller, hvor langt vi er.
        \item Sprint 3 sluttede med 24,5 af 31,5 point. Vi gjorde sprint 4 en uge kort, så vi kunne gøre kortene fra sprint 3 færdige og få rapporten til at passe med koden.
    \end{itemize}
    \item \textbf{Spørgsmål til vejleder}
    \begin{enumerate}
        \item Hvilke diagrammer forventes der i rapporten ud over flow-, sekvens- og pakkediagrammer?
        \item Skal en sprint på én uge bruges på mere kode eller på rapporten?
        \item Hvad synes han, vi skal gå i gang med i sprint 5?
        \item Vi vil gerne kun lave must-have-kravene færdige og droppe should-haves (FR-07 til FR-11 og FR-20) og could-haves. Så kan sprint 6 bruges på accepttest og rapport. Er det i orden, hvis vi begrunder det i diskussionen?
    \end{enumerate}
    \item \textbf{Næste vejledermøde}
    \item \textbf{Eventuelt}
\end{enumerate}

\end{document}
